\documentclass[
  journal=medium,
  manuscript=Qualcomm-Innovation-Fellowship-Proposal,
  year=22-February-2023,
]{cup-journal}

\usepackage{amsmath}
% \usepackage[nopatch]{microtype}
\usepackage{booktabs}
\usepackage{parskip}
\pagenumbering{arabic}

\title{Instrumental Variables for High-Dimensional Data}

\author{Dan Andrei Iliescu (dai24@cam.ac.uk)}
\affiliation{Department of Computer Science and Technology, University of Cambridge, UK}

\addbibresource{example.bib}

% \keywords{keyword entry 1, keyword entry 2, keyword entry 3} %% First letter not capped

\begin{document}

\begin{abstract}
Confounding phenomena negatively affect the prediction quality of many ML models.
One solution to this problem is training a predictive model with an instrumental variable.
This solution is widespread in statistics, economics and epidemiology.
However, instrumental variables are not a widely used method in Machine Learning.
This  
1. An instrumental variable has to be a priori identified by the designer of the experiment. The model itself cannot choose an instrumental variable.
2. The mechanism that enables the identification of the correct causal relationship using an instrumental variable only works for linear models, which are too simple for high-dimensional data.
3. How do we know if the instrumental variable we've chosen is not too restrictive?
In my proposed research, I aim to make the method of instrumental variables widely applicable to solve the problem of confounding in high-dimensional data.
\end{abstract}

\section{Case study: Reinforcement Learning for Human Feedback}

In this case study, we show how a confounder can seriously damage the predictions made by modelling high-dimensional data. 

Consider the task of fine-tuning a language model using human feedback \citep{ouyang2022training}. The goal is to improve how a chatbot answers questions by preforming reinforcement learning on the ratings that users give.

\paragraph{How does it work?} The pre-trained language model is fine-tuned to maximise the ratings it receives after interacting with many human users. Each episode of interaction starts with the user asking a question, the chatbot then provides an answer, and finally the user rates from 0 to 1 how satisfied they are with the answer.

The setup is as follows: the pre-trained language model is fine-tuned to maximise the rewards it receives from a ``preference'' model. The preference model is trained to simulate the ratings users give to the language model's replies.

Each user asks a question (a prompt) and rates on a scale from 0 to 1 the answer they received from the chatbot. Then, a ``preference'' model is trained to predict the user rating conditional on the prompt and the reply.

\paragraph{What's the problem?} A serious problem with this setup is that the preference model learns to discourage certain kinds of replies that are actually helpful. It happens because 

\section{Introduction}

\textbf{Introduction, which describes the problem and current state of the art (i.e. limitations or lack of current solutions)}

There is a disconnect between the problems for which machine learning is designed and the ones on which they are applied.

The problem we are addressing is the classic machine learning task of predicting the target Y from the source X. Alternatively, it is the statistical task of estimating a parameter $\theta$ that maximises the conditional probability of $Y$ given $X$: $P_\theta(Y|X)$.

This task is frequently complicated by the presence of confounders. Confounders are phenomena that distort the relationship between source and target. They arise as random variables that cause both the source and the target.

One solution to this is finding an instrumental variable. This is a variable that causes both the source and target.

However, a major problem with this solution is that instrumental variables have to be independent of the confounder.

\section{Proposal}

\textbf{Proposal, which describes the innovation idea, and what you propose to do.}

I propose to learn an instrumental variable as a latent variable that can be inferred from the environment and then used as conditioning for prediction.

I propose to treat environments as random variables whose values can be inferred and then used as conditioning 

There is a way of training the model that trains for independence instead of maximising the likelihood of the data. It trains a predictive model as a side effect of this independence.

The impact of this innovation would be that one can train a predictive model independent of confounders automatically.

There are 2 goals 

\section{One-Year Horizon}

\textbf{One-year horizon, which explains the target goals for one year and preferably also includes the long-term milestones}

The question is how 

We intend to 

\section{About Me}

\textbf{Description of the student's strength and proposed contribution to the innovation for achieving the milestones. Also describe why this student is better suited for this proposal than anyone else.}


My PhD research explores the interface of generative models and causal inference. Namely I look at how concepts from causal inference help us deal with confounders in domain adaptation.

I also have experience applying deep generative models to real problems in the field of computer vision and speech science.

\printbibliography

\end{document}
